\section{平均场近似用独立性换可计算性}
\label{sec:13-Machine-Learning/13-Statistical-Physics-and-Free-Energy/02}

你给一个层次模型跑平均场变分推断，后验均值和 MCMC 对得上，可信区间却窄得离谱：名义 \(95\%\) 的区间在模拟数据上只覆盖了六成真值。你加大迭代次数、换初始化、调学习率，区间该多窄还是多窄。这不是没收敛，是你挑的那个分布族里压根装不下你要的东西。上一节说变分原理允许把 \(q\) 限制到任意便于计算的族里，本节要看清最常用的那个限制是什么、它换来了什么、又赔进去了什么。

\subsection{平均场就是把联合分布拆成边缘分布的乘积}

多体系统难算的原因只有一条：变量之间有耦合，联合分布不能分解，所以任何期望都要在指数大的状态空间上求和。平均场近似的做法是在上一节的变分问题里加一条约束，只允许乘积形式的分布

\[
q(z)=\prod_{i=1}^{n}q_i(z_i),
\]

其余不作任何假设，每个 \(q_i\) 仍是自由的一元分布。这条约束把``分布''的自由度从指数多个降到 \(n\) 组一元分布，期望和熵都随之分解：\(S(q)=\sum_i S(q_i)\)，而 \(\E_q\) 只需在每个坐标上分别积。

在这个族上极小化变分自由能，可以逐个坐标地做。固定除 \(q_i\) 以外的所有因子，把 \(F[q]\) 中与 \(q_i\) 有关的部分收拢，得到

\[
F[q]=\int q_i(z_i)\,\bar E_i(z_i)\,\mathrm{d}z_i+T\int q_i(z_i)\log q_i(z_i)\,\mathrm{d}z_i+\text{与 }q_i\text{ 无关的项},
\qquad
\bar E_i(z_i)=\E_{q_{-i}}\bigl[E(z)\given z_i\bigr].
\]

括号里剩下的正是一个以 \(\bar E_i\) 为能量的单变量自由能，所以上一节的变分原理直接适用：极小点是这个有效能量的 Boltzmann 分布，

\[
q_i^{*}(z_i)\propto\exp\Bigl(-\frac{1}{T}\E_{q_{-i}}\bigl[E(z)\given z_i\bigr]\Bigr).
\]

取 \(T=1\)、\(E=-\log p(x,z)\)，右边就是变分推断教材里那个坐标上升公式 \(q_i^{*}(z_i)\propto\exp\bigl(\E_{q_{-i}}[\log p(x,z)]\bigr)\)。每个变量只需知道其他变量的``平均效果''就能更新自己，而其他变量的平均效果又依赖于它——这组互相引用的等式就是自洽方程，轮流迭代直到不动。

\subsection{Curie--Weiss 把自洽方程一路算到底}

物理里最干净的例子是 Curie--Weiss 模型：\(N\) 个 Ising 自旋 \(s_i\in\set{-1,+1}\)，两两全连接，耦合按 \(1/N\) 缩放，

\[
E(s)=-\frac{J}{2N}\sum_{i,j}s_is_j-h\sum_i s_i.
\]

取乘积分布 \(q(s)=\prod_i q_i(s_i)\)，其中 \(q_i(s_i)=(1+m_is_i)/2\)，于是 \(m_i=\E_{q_i}[s_i]\) 是唯一的自由参数。代入变分自由能，能量项因独立性化为 \(-\frac{J}{2N}\sum_{i,j}m_im_j-h\sum_i m_i\)（对角项 \(s_i^2=1\) 只贡献常数），熵项化为各自的二元熵之和。对 \(m_i\) 求导置零，并利用对称性设所有 \(m_i=m\)，得到

\[
m=\tanh\Bigl(\frac{Jm+h}{T}\Bigr).
\]

这就是平均场方程。它的读法是：每个自旋不再逐一感受其他自旋，而是感受一个有效场 \(Jm+h\)，场的强度又由整体磁化 \(m\) 决定。方程右边的 \(\tanh\) 来自二元熵的导数，左边的 \(m\) 来自能量项，温度 \(T\) 站在两者之间当兑换率——和上一节那张图右半边是同一件事。

\(h=0\) 时方程退化为 \(m=\tanh(Jm/T)\)。\(T\ge J\) 时只有 \(m=0\) 一个解，自旋各行其是；\(T<J\) 时零解失稳，出现一对对称的非零解 \(\pm m^{*}\)，系统整体倒向一边。值得强调的是，在这个全连接模型的 \(N\to\infty\) 极限下，平均场给出的自由能密度与自洽方程是精确的，因为每个自旋连着 \(O(N)\) 个邻居而每根耦合只有 \(O(1/N)\) 强，邻居之和的相对涨落被平均掉了。平均场不是普遍成立的定理，它是这一类极限下的精确结论、别处的一阶近似。

\subsection{独立假设的账单是被系统性低估的涨落}

乘积形式换来了可计算性，要付的写在方差上，而且方向是固定的：偏小，不是随机偏。高斯情形可以把这句话证成一条命题。

\begin{proposition}
设 \(p\) 是均值任意、精度矩阵为 \(\Lambda=\Sigma^{-1}\) 的非退化高斯分布。在乘积族中极小化 \(\KL(q\|p)\) 的解 \(q^{*}=\prod_i q_i^{*}\) 满足 \(\Var_{q_i^{*}}(z_i)=1/\Lambda_{ii}\)，因而对每个 \(i\) 都有 \(\Var_{q_i^{*}}(z_i)\le\Sigma_{ii}\)，等号成立当且仅当 \(\Lambda_{i,-i}=0\)，即 \(z_i\) 与其余变量在 \(p\) 下独立。
\end{proposition}

骨架只有两步。第一步，把上一小节的坐标更新用到高斯上：\(\log p\) 是 \(z\) 的二次型，\(\E_{q_{-i}}[\log p]\) 关于 \(z_i\) 仍是二次的，二次项系数是 \(-\Lambda_{ii}/2\) 且不含其他因子的信息，所以 \(q_i^{*}\) 是方差 \(1/\Lambda_{ii}\) 的高斯。第二步，用分块求逆的 Schur 补公式

\[
\Sigma_{ii}=\bigl(\Lambda_{ii}-\Lambda_{i,-i}\Lambda_{-i,-i}^{-1}\Lambda_{i,-i}^{\trans}\bigr)^{-1},
\]

括号里减掉的那一项因 \(\Lambda_{-i,-i}\succ0\) 而非负，于是 \(\Sigma_{ii}\ge1/\Lambda_{ii}\)。条件对账：需要 \(\Sigma\) 正定（否则 \(\Lambda\) 不存在）、\(\Lambda_{-i,-i}\) 可逆（由 \(\Lambda\succ0\) 保证）、以及乘积族里的极小确实取到（高斯情形下坐标更新是线性方程组，解存在唯一）。三条都是可查的，不是``一般来说''。

数字上感受一下。取二维、相关系数 \(\rho=0.9\)、两个边缘方差都是 \(1\)，则 \(\Lambda_{ii}=1/(1-\rho^2)=5.26\)，平均场给出的边缘标准差是 \(\sqrt{1-\rho^2}=0.44\)。真值是 \(1\)。你的可信区间窄了两倍多，覆盖率掉到六成就不奇怪了。

\begin{center}
\begin{tikzpicture}[>=stealth,line join=round]
% ---- 左：相关高斯，平均场收窄 ----
\draw[gray,thin] (-1.7,0) -- (1.7,0);
\draw[gray,thin] (0,-1.25) -- (0,1.25);
\draw[thick,dashed,rotate=45] (0,0) ellipse (1.378 and 0.316);
\draw[thick] (0,0) circle (0.436);
\node[anchor=north east,font=\small] at (-0.95,-0.95) {真后验 \(p\)};
\draw[->] (1.15,-0.85) node[anchor=west,font=\small] {平均场 \(q\)} -- (0.32,-0.32);
\draw[thick,dashed] (-1,-1.62) -- (1,-1.62);
\node[anchor=west,font=\small] at (1.06,-1.62) {\(p\)};
\draw[thick] (-0.436,-1.94) -- (0.436,-1.94);
\node[anchor=west,font=\small] at (0.50,-1.94) {\(q\)};
\node[anchor=north,font=\small] at (0,-2.12) {边缘标准差 \(1\) 被压成 \(0.44\)};
% ---- 右：双峰后验，平均场只抱住一个峰 ----
\begin{scope}[shift={(5.2cm,-0.95cm)},x=0.8cm]
  \draw[->] (-3.3,0) -- (3.4,0) node[below,font=\small] {\(z\)};
  \draw[thick,dashed] plot[smooth] coordinates {
    (-3.0,0.011) (-2.6,0.066) (-2.2,0.249) (-1.8,0.605) (-1.4,0.943)
    (-1.2,0.998) (-1.0,0.944) (-0.6,0.616) (-0.2,0.314) (0,0.270)
    (0.2,0.314) (0.6,0.616) (1.0,0.944) (1.2,0.998) (1.4,0.943)
    (1.8,0.605) (2.2,0.249) (2.6,0.066) (3.0,0.011)};
  \draw[thick] plot[smooth] coordinates {
    (-3.0,0.022) (-2.6,0.131) (-2.2,0.497) (-1.8,1.210) (-1.4,1.887)
    (-1.2,1.994) (-1.0,1.887) (-0.6,1.210) (-0.2,0.497) (0,0.270)
    (0.2,0.131) (0.6,0.022) (1.0,0.002) (1.6,0.000) (2.4,0.000)
    (3.0,0.000)};
  \node[anchor=west,font=\small] at (-3.3,2.22) {平均场解 \(q\)};
  \node[anchor=west,font=\small] at (1.45,1.05) {真后验 \(p\)};
\end{scope}
\end{tikzpicture}
\end{center}

\emph{左：相关被拆掉之后，等高线从斜椭圆缩成坐标轴对齐的小圆；右：多峰时同一个准则挑走一个峰，另一个峰的质量被整个丢弃。}

\subsection{低估涨落与反向 KL 的模式寻找是同一句话}

物理文献说``平均场低估涨落''，机器学习文献说``反向 KL 是 mode-seeking 的''，这两句话描述的是同一个机制。上一节已经把变分自由能写成 \(F[q]=T\KL(q\|p)+F\)，所以极小化变分自由能就是极小化 \(\KL(q\|p)\)，KL 的第一个位置放的是近似分布。展开这个方向的 KL，被积的权重是 \(q\)：凡是 \(q\) 取正而 \(p\) 近乎为零的地方，\(\log(q/p)\) 爆掉，罚得极重；反过来，\(p\) 有质量而 \(q\) 为零的地方，权重 \(q=0\)，一分钱不罚。方向的这种不对称参见\hyperref[sec:07-Information-Theory/03-Divergences/02]{``Gibbs 不等式与 KL 的方向性''}。

于是最优的 \(q\) 只敢待在 \(p\) 确实有质量的区域里，宁可漏掉也不肯溢出。这条行为准则在两种几何下呈现为两副面孔。当 \(p\) 是单峰但强相关的，它的高密度区是一条细斜带，而乘积族只能画坐标轴对齐的等高线；要塞进斜带里就必须两个方向一起缩，缩出来的正是那个 \(1/\Lambda_{ii}\)。当 \(p\) 是多峰的，两峰之间是低密度谷，横跨两峰的 \(q\) 会在谷里挨罚，于是最优解干脆抱住一个峰。图的左右两半画的就是这两副面孔：一个叫低估涨落，一个叫挑模式，源头是同一个不对称的散度。

这个统一有实际用处。它告诉你后验区间偏窄不是数值问题，加迭代、换优化器都救不了，得改分布族——引入结构化 \(q\)（保留部分依赖）、用归一化流、或者改用采样。它也告诉你 \(q\) 报出的不确定性不能直接当作真实不确定性使用，尤其在下游要做决策的时候。变分推断这一侧的完整讨论见\hyperref[sec:13-Machine-Learning/08-Sampling-and-Inference/03]{``变分推断把后验近似变成优化''}。

\subsection{平均场在机器学习里落在哪些地方}

受限玻尔兹曼机是最直接的例子。它的二分图结构使得给定可见层时隐层各单元条件独立，这一半不需要近似；难的是训练梯度里的负相，也就是上一节说的自由能的导数，它要模型自身分布下的期望。平均场用乘积分布近似这个分布，持续对比散度改用一条短 MCMC 链，两者是估同一项的两条路，前者偏差固定、方差为零，后者偏差随链长下降、方差不为零。

主题模型、隐 Markov 结构、稀疏编码里的变分 E 步，也都在跑同一套坐标更新。写代码时它们长得像 EM：轮流更新每个因子，每轮保证 ELBO 不降，所以可以直接用 ELBO 的曲线判断是否收敛。要注意的是，ELBO 收敛只说明你在这个族里找到了（局部）最优，不说明这个族够用——覆盖率那类校准检查是另一回事，必须单独做。

平均场的价值在于它是第一个非平凡近似，而且在全连接、高温这类极限下是精确的。真正需要小心的是它失效的方式：不是均匀地差一点，而是在某些参数区域里定性地错。这些区域在哪里、错成什么样，是下一节的内容。
